
\section{Related Work}
\label{sec:related_work}
\seckiny

Building hierarchical agents is a long standing topic in reinforcement learning~\cite{sutton1999between,precup2000temporal, dayan1993feudal, dietterich2000hierarchical, boutilier1997prioritized, dayan1993improving, kaelbling2014hierarchical, parr1998reinforcement, precup1997planning, precup1998theoretical, schmidhuber1991neural, sutton1995td, wiering1997hq,vezhnevets2016straw,bacon2015option}. 
The options framework~\cite{sutton1999between,precup2000temporal} is a popular formulation for considering the problem with a two level hierarchy. The bottom level -- an option -- is a sub-policy with a termination condition, which takes in environment observations and outputs actions until the termination condition is met. 
An agent picks an option using its policy-over-options (the top level) and subsequently follows it until termination, at which point the policy-over-options is queried again and the process continues. 
Options are typically learned using sub-goals and `pseudo-rewards' that are provided explicitly~\cite{sutton1999between, dietterich2000hierarchical, dayan1993feudal}. For a simple, tabular case~\cite{wiering1997hq,schaul2015uvfa}, each state can be used as a sub-goal.
Given the options, a policy-over-options can be learned using standard techniques by treating options as actions. 
Recently~\cite{tessler2016minecraft, tejas2016hdrl} have demonstrated that combining deep learning with pre-defined sub-goals delivers promising results in challenging environments like Minecraft and Atari, however sub-goal discovery was not addressed.

A recent work of~\cite{bacon2015option} shows the possibility of learning options jointly with a policy-over-options in an end-to-end fashion by extending the policy gradient theorem to options. 
When options are learnt end-to-end, they tend to degenerate to one of two trivial solutions: (i) only one active option that solves the whole task; (ii) a policy-over-options that changes options at every step, micro-managing the behaviour. Consequently, regularisers~\cite{bacon2015option,vezhnevets2016straw} are usually introduced to steer the solution towards multiple options of extended length. This is believed to provide an inductive bias towards re-usable temporal abstractions and to help generalisation. 

A key difference between our approach and the options framework is that in our proposal the top level produces a meaningful and explicit goal for the bottom level to achieve. Sub-goals emerge as directions in the latent state-space and are naturally diverse. We also achieve significantly better scores on ATARI than Option-Critic (section~\ref{sec:Experiments}).


There has also been a significant progress in non-hierarchical deep RL methods by using auxiliary losses and rewards. \cite{bellemare2016countexplore} have significantly advanced the state-of-the-art on Montezuma's Revenge by using pseudo-count based auxiliary rewards for exploration, which stimulate agents to explore new parts of the state space. The recently proposed UNREAL agent~\cite{jaderberg2016unreal} also demonstrates a strong improvement by using unsupervised auxiliary tasks to help refine its internal representations. We note that these benefits are orthogonal to those provided by FuN, and that both approaches could be combined with FuN for even greater effect. 

